\setcounter{section}{2}

\Needspace{5\baselineskip}
\hypertarget{alphaevolve}{%
\section{AlphaEvolve}\label{alphaevolve}}

\Needspace{5\baselineskip}
\hypertarget{i.-the-core-idea-2}{%
\subsection{I. The Core Idea}\label{i.-the-core-idea-2}}

While some large-model approaches combine LLMs with search algorithms, AlphaEvolve combines them with optimization resembling genetic algorithms. Using an existing codebase as context, it generates many programs with an LLM, evaluates them with a verifier, and stores better-scoring code in the codebase, continually improving through iterations within the codebase (context).

Rather than relying on manually designed mutation operators as traditional genetic algorithms do, AlphaEvolve uses the LLM's world knowledge and programming capabilities as an ``intelligent mutation operator'' that makes meaningful code modifications. It can also be regarded as policy improvement through context updates while keeping model weights unchanged.

\Needspace{5\baselineskip}
\hypertarget{ii.-workflow}{%
\subsection{II. Workflow}\label{ii.-workflow}}

\begin{enumerate}
\def\labelenumi{\arabic{enumi}.}
\item
  Initialization: human scientists define the problem, provide initial code, and mark the code blocks to evolve. A dynamic database is initialized to simulate a population and store different programs. An automatic evaluation function scores generated code.
\item
  Sampling and prompt construction: sample a ``parent program'' and several ``inspiration programs'' from the database. The prompt sampler combines these, past excellent solutions, and system instructions into a context-rich prompt for the LLM.
\item
  LLM generation and mutation: propose changes to the parent program according to the prompt. Outputs are usually diffs, containing only code changes, which an automated script converts into a complete modified child-program file.
\item
  Code execution and evaluation: send the new program to an evaluator for actual execution. A user-defined evaluate function calculates one or more metrics, such as accuracy, speed, or mathematical properties. For time-consuming tasks, the system supports cascaded evaluation: quickly filter with small-scale tests before running large-scale tests.
\item
  Database update: if the new program performs well, store it in the program database.
\end{enumerate}

\Needspace{5\baselineskip}
\hypertarget{iii.-applications}{%
\subsection{III. Applications}\label{iii.-applications}}

Code is the example above, but as an optimization algorithm whose operator is simply replaced by an LLM, it can evolve objects at different levels: directly optimize strings or numerical values, write a function that generates solutions, or evolve an algorithm for searching for answers. Evolving algorithms as ``instrumental objectives'' is often more effective than direct search.

AlphaEvolve can also evolve prompts themselves. It treats LLM system instructions and context as optimization targets, continually improving ``how to ask the LLM'' in a separate parallel evolutionary process to obtain better code generation.

\Needspace{5\baselineskip}
\hypertarget{iv.-key-components-and-strategies}{%
\subsection{IV. Key Components and Strategies}\label{iv.-key-components-and-strategies}}

\Needspace{4\baselineskip}
\begin{enumerate}
\def\labelenumi{(\Roman{enumi})}
\tightlist
\item
  LLM ensemble strategy
\end{enumerate}

Flash: fast and low-latency, used to generate many candidate programs and rapidly explore the search space.

Pro: more capable, used less frequently to provide high-quality, deep code modifications that help evolution break through bottlenecks and achieve ``qualitative leaps.''

Code generated by both Flash and Pro is sent to evaluators, with retention determined by objective scores.

\Needspace{4\baselineskip}
\begin{enumerate}
\def\labelenumi{(\Roman{enumi})}
\setcounter{enumi}{1}
\tightlist
\item
  Asynchronous concurrency
\end{enumerate}

The system implements asynchronous concurrency using asyncio. The controller, LLM samplers, and evaluators are decoupled to maximize throughput. Even if one evaluation takes hours, such as model training, it does not block the entire evolutionary process.

\Needspace{4\baselineskip}
\begin{enumerate}
\def\labelenumi{(\Roman{enumi})}
\setcounter{enumi}{2}
\tightlist
\item
  Database maintenance
\end{enumerate}

Use MAP-Elites-like algorithms to divide the solution space into ``niches.''

\Needspace{4\baselineskip}
\begin{enumerate}
\def\labelenumi{\arabic{enumi}.}
\tightlist
\item
  Defining niches through multiple evaluation metrics
\end{enumerate}

\begin{enumerate}
\def\labelenumi{(\arabic{enumi})}
\tightlist
\item
  User-provided evaluate function: returns a dictionary of scalar values, such as accuracy, runtime speed, memory consumption, and code length. These form a feature space with multiple optimizable metrics; different combinations define niches. For example, one niche may store ``extremely fast but lengthy code,'' another ``slightly slower code with extremely low memory consumption.''
\end{enumerate}

AlphaEvolve explicitly notes that optimizing multiple metrics, even when only one matters, helps discover programs with different structures or logic. In other words, \textbf{this algorithm can optimize multiple variables.}

\begin{enumerate}
\def\labelenumi{(\arabic{enumi})}
\setcounter{enumi}{1}
\tightlist
\item
  LLM scoring: an LLM can score code conciseness, readability, or other stylistic features. These scores are added to the evaluation dictionary as extra niche-defining dimensions. For example, code with ``very clear and concise logic'' may be retained despite slightly weaker performance because it occupies a ``high-readability'' niche.
\end{enumerate}

\Needspace{4\baselineskip}
\begin{enumerate}
\def\labelenumi{\arabic{enumi}.}
\setcounter{enumi}{1}
\tightlist
\item
  Principles for retaining solutions each round
\end{enumerate}

Exploitation: when a new solution outperforms the existing solution in a particular niche, it replaces it; otherwise, it is discarded. In short, retain the best-performing solution.

Exploration: retain structurally distinctive solutions occupying niches different from all existing solutions to ensure database diversity.

The database's main purpose is to optimally reproduce previously explored ideas for use in future generations. It therefore favors valuable, high-quality, or distinctive ``elite'' code while preserving as much diversity as possible to avoid premature convergence to local optima, rather than storing every historically generated piece of poor code indefinitely.

\Needspace{4\baselineskip}
\begin{enumerate}
\def\labelenumi{\arabic{enumi}.}
\setcounter{enumi}{2}
\tightlist
\item
  Physical isolation through an island model
\end{enumerate}

In addition to feature-based niches, AlphaEvolve incorporates an island model, physically dividing the population into separate ``islands'' (subpopulations) to create isolated ecological environments. Even programs with similar metrics can coexist without eliminating each other if they belong to different islands. This enforces diversity of genetic origins.

\Needspace{4\baselineskip}
\begin{enumerate}
\def\labelenumi{(\Roman{enumi})}
\setcounter{enumi}{3}
\tightlist
\item
  Selecting parent and inspiration programs
\end{enumerate}

\Needspace{4\baselineskip}
\begin{enumerate}
\def\labelenumi{\arabic{enumi}.}
\tightlist
\item
  Parent-program selection
\end{enumerate}

The parent is the base code directly modified by the LLM. Selection balances exploitation and exploration, with relatively more exploitation because inspirations also provide an exploration mechanism.

Exploitation: frequently sample the highest-scoring, best-performing programs and let the LLM refine them to push scores higher.

Exploration: to avoid local optima, also sample programs with ``acceptable performance but distinctive structure'' or those occupying different niches.

Island-model rotation: if an island model is used, parents may come from different islands, that is, independent evolving subpopulations. This maintains genetic diversity and prevents all parents from converging.

\Needspace{4\baselineskip}
\begin{enumerate}
\def\labelenumi{\arabic{enumi}.}
\setcounter{enumi}{1}
\tightlist
\item
  Inspiration-program selection
\end{enumerate}

Inspiration programs are shown to the LLM as contextual examples in the prompt. They are references, not directly modified. Diversity is a key selection principle. Particularly in multi-objective optimization, such as jointly optimizing speed and memory, the system selects programs excelling on different metrics. For example, an extremely fast but lengthy program and a very concise but slightly slower one are placed together in the prompt, inspiring a novel solution combining their strengths or providing ``inspiration'' through their differences.

\Needspace{4\baselineskip}
\begin{enumerate}
\def\labelenumi{\arabic{enumi}.}
\setcounter{enumi}{2}
\tightlist
\item
  Technical background
\end{enumerate}

AlphaEvolve does not merely sample old code randomly; its core combines in-context learning with evolutionary strategies. The parent determines the evolutionary starting point, and inspirations guide the direction. By carefully selecting inspirations, the system constructs a dynamic prompt, using the LLM's strong contextual understanding to explicitly convey implicit knowledge in the database, such as a programming pattern effective for the current task, thereby increasing the probability of generating high-quality diffs.

\FloatBarrier
\Needspace{5\baselineskip}
\hypertarget{references-121}{%
\subsection{References}\label{references-121}}

\vspace{0.25em}
\begin{itemize}
\tightlist
\item
  Novikov, A., Vu, N., Eisenberger, M., et al.~(2025). \href{https://arxiv.org/abs/2506.13131}{AlphaEvolve: A coding agent for scientific and algorithmic discovery}. arXiv:2506.13131.
\end{itemize}

\clearpage
